% ============================================================================
%  Chapitre 6 — Évaluation et résultats
%  Cible : 10 à 12 pages. À rédiger J10-J11.
%
%  AVERTISSEMENT : ce chapitre ne peut être rédigé qu'une fois les mesures
%  collectées (voir scripts/collecte-mesures.sh et annexe B).
%  Aucune valeur ne doit être inventée. Si une mesure « avant » n'est pas
%  disponible, l'écrire explicitement et en faire une limite méthodologique.
% ============================================================================
\chapter{Évaluation et résultats}
\label{chap:evaluation}

\section{Protocole d'évaluation}
\label{sec:eval-protocole}

\subsection{Questions d'évaluation}
L'évaluation vise à répondre aux quatre sous-questions posées en
section~\ref{sec:intro-problematique} :

\begin{description}[style=nextline,leftmargin=1.5cm]
  \item[QE1] L'environnement de construction obtenu est-il effectivement
    reproductible ?
  \item[QE2] L'intégration des analyses de sécurité améliore-t-elle la
    détection des défauts sans dégrader la vélocité ?
  \item[QE3] Les indicateurs collectés permettent-ils de piloter la qualité ?
  \item[QE4] L'infrastructure d'exécution est-elle maîtrisée en coût et en
    effort d'administration ?
\end{description}

\subsection{Sources de mesure}
\begin{table}[H]
\centering
\caption{Sources des mesures utilisées}
\label{tab:eval-sources}
\begin{tabularx}{\textwidth}{@{}l X l@{}}
\toprule
\textbf{Indicateur} & \textbf{Source} & \textbf{Période} \\
\midrule
Durée et issue des exécutions & Interface de la plateforme d'orchestration & \\
Couverture de code & Rapports de l'outil de couverture & \\
Défauts d'analyse statique & Vues du serveur d'analyse & \\
Vulnérabilités de composants & Rapports de l'analyseur binaire & \\
Effort d'administration & Relevé déclaratif auprès de l'équipe & \\
\bottomrule
\end{tabularx}
\end{table}

\aecrire{Préciser la période d'observation, le nombre d'exécutions prises en
compte, et le traitement des valeurs aberrantes (exécutions interrompues,
indisponibilité d'un service externe). Une évaluation sans protocole explicite
n'est pas défendable.}

\subsection{Validité}
\aecrire{Section courte mais indispensable. Traiter : validité interne (les
améliorations sont-elles attribuables à la chaîne ou à d'autres évolutions
concomitantes ?), validité externe (généralisation à d'autres produits),
fiabilité (les mesures sont-elles reproductibles ?).}

% ===========================================================================
\section{Impact}
\label{sec:eval-impact}

\aecrire{Section chapeau (2-3 phrases) : annoncer que cette section restitue
les effets de la chaîne au-delà de la seule mesure technique — indicateurs de
performance de livraison, coût, retour des équipes.}

\subsection{Indicateurs de performance de livraison}
\label{sec:eval-dora}

\begin{table}[H]
\centering
\caption{Indicateurs de performance adaptés au contexte embarqué}
\label{tab:eval-dora}
\begin{tabularx}{\textwidth}{@{}l X c c@{}}
\toprule
\textbf{Indicateur} & \textbf{Adaptation retenue} & \textbf{Avant} & \textbf{Après} \\
\midrule
Fréquence d'intégration & Nombre d'intégrations validées par semaine & & \\
Délai de retour & Délai entre soumission et résultat complet & & \\
Taux d'échec & Part des exécutions en échec imputables au code & & \\
Délai de rétablissement & Temps de remise en état d'une chaîne cassée & & \\
\bottomrule
\end{tabularx}
\end{table}

\aecrire{Justifier chaque adaptation par rapport à la définition d'origine
(section~\ref{sec:art-dora}). Si les valeurs « avant » sont estimées et non
mesurées, l'indiquer dans le tableau et dans le texte.}

\subsection{Coût et effort d'administration}
\label{sec:eval-cout}

\aecrire{Estimer : coût d'exécution de l'infrastructure, temps d'ingénierie
économisé sur les tâches auparavant manuelles (analyse de composants avant
jalon, préparation des environnements, collecte des rapports), et effort de
maintenance de la chaîne elle-même. Ne pas dissimuler ce dernier poste : une
chaîne d'intégration a un coût de possession.}

\subsection{Retour des utilisateurs}
\label{sec:eval-retour}

\aecrire{Si un recueil auprès des développeurs est possible dans les délais,
même informel, l'intégrer : quelques verbatims bien choisis renforcent
considérablement une étude de cas. Préciser le nombre de personnes
interrogées et la méthode.}

\subsection{Apports pour l'entreprise}
\label{sec:eval-apports-entreprise}

\aecrire{Distinguer l'apport technique de l'apport organisationnel : au-delà
de l'outillage, la chaîne modifie les pratiques de l'équipe, rend visible la
qualité et fournit les éléments de preuve attendus par le cadre normatif.}

% ===========================================================================
\section{Difficultés rencontrées}
\label{sec:eval-difficultes}

\aecrire{Section chapeau : consolider ici les difficultés concrètes
rencontrées pendant la mise en œuvre (renvoyer aux difficultés déjà détaillées
section~\ref{sec:impl-pipeline-app} pour la chaîne applicative) et celles
propres à la mesure elle-même.}

\subsection{Difficultés techniques}
\aecrire{Reprendre et synthétiser les difficultés déjà documentées au
chapitre~5 : nettoyage des services en fin d'exécution sur un exécuteur neuf,
délai d'attente au démarrage des services, exclusion des tests destinés au
cœur temps réel, normalisation des chemins pour l'analyse statique sur
exécuteurs jetables.}

\subsection{Difficultés liées à la mesure}
\aecrire{Distinguer des limites méthodologiques (traitées en section
Discussion) : il s'agit ici des obstacles concrets rencontrés pour collecter
les mesures elles-mêmes — étude de cas unique, mesures partiellement
déclaratives, effet d'apprentissage de l'équipe concomitant à la période
d'observation.}

% ===========================================================================
\section{Résultats}
\label{sec:eval-resultats}

\aecrire{Section chapeau (2-3 phrases) : annoncer que cette section présente
les résultats quantitatifs obtenus, dans l'ordre reproductibilité,
performance, puis qualité et sécurité.}

\subsection{Reproductibilité de l'environnement}
\label{sec:eval-reproductibilite}

\aecrire{Protocole proposé : construire une même révision du code à deux dates
distinctes et sur deux machines distinctes avec la même étiquette d'image,
puis comparer. Indiquer précisément ce qui est comparé : versions des outils,
réussite de la construction, et si possible empreinte des artefacts produits.
Discuter les écarts constatés et leurs causes (horodatages, chemins de
construction, ordre de liaison).}

\chiffre{Résultat de la comparaison des empreintes d'artefacts.}

\subsection{Performance de la chaîne}
\label{sec:eval-performance}

\subsubsection{Durées de construction}
\begin{table}[H]
\centering
\caption{Durées de construction observées}
\label{tab:eval-durees}
\begin{tabular}{@{}lccc@{}}
\toprule
\textbf{Chaîne} & \textbf{Minimum} & \textbf{Médiane} & \textbf{Maximum} \\
\midrule
Applicative (matrice complète) & & & \\
Distribution, variante dev & & & \\
Distribution, variante prod & & & \\
Publication de l'image de compilation & & & \\
\bottomrule
\end{tabular}
\end{table}

\subsubsection{Effet du cache d'état partagé}
\aecrire{Comparer construction à froid et à chaud. C'est un résultat
quantitatif fort et facile à obtenir : ne pas le négliger. Présenter sous
forme de tableau puis d'un graphique.}

\begin{figure}[H]
\centering
\aecrire{Graphique : durée de construction avec et sans cache, par variante.}
\caption{Incidence du cache d'état partagé sur la durée de construction}
\label{fig:eval-cache}
\end{figure}

\subsubsection{Effet du parallélisme de la matrice}
\aecrire{Comparer la durée cumulée des quatre configurations exécutées en
séquence à la durée réelle observée en parallèle.}

\subsection{Qualité et sécurité}
\label{sec:eval-qualite}

\subsubsection{Couverture de code}
\begin{figure}[H]
\centering
\aecrire{Graphique d'évolution de la couverture (lignes, fonctions, branches)
sur la période d'observation.}
\caption{Évolution de la couverture de code}
\label{fig:eval-couverture}
\end{figure}

\subsubsection{Défauts détectés par l'analyse statique}
\aecrire{Présenter la répartition initiale des défauts par priorité, puis leur
évolution. Distinguer les défauts préexistants des défauts nouvellement
introduits : c'est la distinction qui donne du sens à la mesure.}

\chiffre{Nombre de défauts de priorité 1 et 2 au démarrage et en fin de
période.}

\subsubsection{Vulnérabilités des composants tiers}
\aecrire{Présenter le nombre de composants identifiés et la répartition des
vulnérabilités par sévérité. Souligner l'apport principal : avant la mise en
place, l'inventaire des composants n'existait pas de manière automatisée.}

\subsubsection{Détection des régressions}
\aecrire{Si des cas concrets de régressions détectées par la chaîne avant
intégration sont documentés, les citer. Un exemple réel vaut mieux qu'un
paragraphe général.}

% ===========================================================================
\section{Discussion}
\label{sec:eval-discussion}

\subsection{Réponses aux questions d'évaluation}
\aecrire{Reprendre QE1 à QE4 et y répondre explicitement, en s'appuyant sur
les mesures présentées. Une réponse nuancée et argumentée est attendue.}


\subsection{Limites}
\aecrire{Recenser honnêtement : absence de tests sur cible matérielle réelle
dans la chaîne, politique de blocage non appliquée sur l'analyse statique,
exécution en superutilisateur, absence de mesure de référence antérieure pour
certains indicateurs, période d'observation courte. Une section de limites
solide est un marqueur de maturité, elle ne dessert pas la note.}

\subsection{Menaces à la validité}
\aecrire{Distinguer des limites : il s'agit ici de ce qui pourrait invalider
les conclusions. Étude de cas unique, mesures partiellement déclaratives,
effet d'apprentissage de l'équipe concomitant.}
